% from ega0-19-fr.tex
\oldpage[0\textsubscript{IV-1}]{71}

Si $M$ et $N$ sont deux $A$-modules topologiques dont les topologies sont toutes
deux déduites de celle de $A$, alors \emph{tout} $A$-homomorphisme $u:M\to N$
est \emph{continu}, car pour tout voisinage $V$ de 0 dans $N$, il existe par
hypothèse un idéal ouvert $\mathfrak J$ de $A$ tel que $\mathfrak JN\subset V$,
et l'on a donc $u(\mathfrak JM)\subset\mathfrak JN\subset V$.

Lorsque $B$ est une $A$-algèbre topologique, la topologie sur $B$ \emph{déduite}
de celle de $A$ est \emph{plus fine} que la topologie donnée, car pour tout
idéal ouvert $\mathfrak R$ de $B$, il y a par hypothèse un idéal ouvert
$\mathfrak J$ de $A$ tel que $\mathfrak JB\subset\mathfrak R$.
\end{env}

\subsection{Épimorphismes et monomorphismes formels}
\label{subsection:0.19.1-fr}

\begin{proposition}[19.1.1]
\label{0.19.1.1-fr}
Soient $A$ un anneau topologique, $M$, $N$ deux $A$-modules topologiques,
$(W_\lambda)$ un système fondamental de sous-modules ouverts dans $N$,
$u:M\to N$ un $A$-homomorphisme continu, $\widehat M$, $\widehat N$ les
séparés complétés de $M$ et $N$, $\widehat u:\widehat M\to\widehat N$ le
prolongement continu de $u$ aux séparés complétés.
\begin{enumerate}
  \item[(i)] Les conditions suivantes sont équivalentes~:
  \begin{enumerate}
    \item[a)] $u(M)$ est \emph{dense} dans $N$.
    \item[a$'$)] $\widehat u(\widehat M)$ est \emph{dense} dans $\widehat N$.
    \item[a$''$)] Pour tout $\lambda$, l'homomorphisme composé
    $M\xrightarrow{u}N\to N/W_\lambda$ est \emph{surjectif}.
  \end{enumerate}
  \item[(ii)] Les conditions suivantes sont équivalentes~:
  \begin{enumerate}
    \item[b)] L'image réciproque par $u$ de la topologie de $N$ est égale à la
    topologie de $M$.
    \item[b$'$)] $\widehat u$ est un \emph{isomorphisme} du $\widehat A$-module
    topologique $\widehat M$ sur le sous-$\widehat A$-module topologique
    $\widehat u(\widehat M)$ de $\widehat N$ (qui est nécessairement
    \emph{fermé}).
    \item[b$''$)] Les $u^{-1}(W_\lambda)$ forment un système fondamental de
    voisinages de 0 dans $M$.
  \end{enumerate}
\end{enumerate}
Cela résulte immédiatement de la définition des séparés complétés, et du fait
que $N/W_\lambda$ est discret.
\end{proposition}

\begin{env}[19.1.2]
\label{0.19.1.2-fr}
Lorsque les conditions équivalentes de (i) (resp. (ii)) dans (19.1.1) sont
remplies, on dit que $u$ est un \emph{épimorphisme formel} (resp. un
\emph{monomorphisme formel}). On dit que $u$ est un \emph{bimorphisme formel}
s'il est à la fois un monomorphisme formel et un épimorphisme formel~; il
revient au même, en vertu de (19.1.1), de dire que $\widehat u$ est un
\emph{isomorphisme} du $\widehat A$-module topologique $\widehat M$ sur le
$\widehat A$-module topologique $\widehat N$.
\end{env}

\begin{proposition}[19.1.3]
\label{0.19.1.3-fr}
Soient $A$ un anneau topologique, $M$, $N$ deux $A$-modules topologiques,
$u:M\to N$ un $A$-homomorphisme continu. On suppose qu'il existe deux systèmes
fondamentaux de sous-modules ouverts, $(V_\lambda)$, $(W_\lambda)$ dans $M$ et
$N$ respectivement, ayant même ensemble d'indices et tels que
$u(V_\lambda)\subset W_\lambda$ pour tout $\lambda$~; soit
$u_\lambda:M/V_\lambda\to N/W_\lambda$ l'homomorphisme déduit de $u$ par
passage aux quotients. Alors~:
\begin{enumerate}
  \item[(i)] Pour que $u$ soit un épimorphisme formel, il faut et il suffit que,
  pour tout $\lambda$, $u_\lambda$ soit surjectif.
  \item[(ii)] Pour que $u$ soit un monomorphisme formel, il faut et il suffit
  que, pour tout $\lambda$, il existe un $\mu$ tel que
  $V_\mu\subset V_\lambda$ et que $\operatorname{Ker}(u_\mu)$ soit contenu
  dans le noyau $V_\lambda/V_\mu$ de l'application canonique
  $M/V_\mu\to M/V_\lambda$.
  \item[(iii)] Pour que $u$ soit un bimorphisme formel, il faut et il suffit
  que, pour tout $\lambda$, $u_\lambda$ soit surjectif et qu'il existe un
  $\mu$ tel que $V_\mu\subset V_\lambda$ et que l'homomorphisme canonique
  $M/V_\mu\to M/V_\lambda$

  